\documentclass[conference]{IEEEtran}
\usepackage{cite}
\usepackage{amsmath,amssymb}
\usepackage{graphicx}
\usepackage{booktabs}
\usepackage{multirow}
\usepackage{siunitx}
\usepackage{hyperref}
\usepackage{xcolor}
\usepackage[T1]{fontenc}
\usepackage{placeins}  % For \FloatBarrier to control float placement

% Enable UTF-8
\usepackage[utf8]{inputenc}
% Restore macro shims to neutralize corrupted table tokens (will allow compilation)
% Removed temporary token-recovery shims (\oprule, \extbf) after table cleanup.

% (Removed temporary corrupted table macros after surgical cleanup)

% Graphics path helper (compiled from paper/)
% Figures live under ../figures/...
\graphicspath{{../figures/}}

% Helper to include a figure if it exists, else show a placeholder box
\makeatletter
% Simplified, robust figure/placeholder macro (avoid compound numeric factors causing 'Illegal unit of measure')
\newcommand{\figorplaceholder}[2][\linewidth]{%% opt arg = width, mandatory arg = file path
  \IfFileExists{../figures/#2}{%% file exists
    \includegraphics[width=#1]{#2}%
  }{%% placeholder
    \fbox{\parbox[c]{#1}{\centering \small Placeholder (figure not found)}}%
  }%% end IfFileExists
}
\makeatother

\title{Predictive Design of Perovskite Materials for Solar Cell Applications}

\author{
\IEEEauthorblockN{R. Chaurasiya, A. Soni, and S. Awasthi}
\IEEEauthorblockA{
Department of Electronics and Communication Engineering\\
Indian Institute of Information Technology, Bhopal, India\\
Emails: rchaurasiya@iiitbhopal.ac.in, aishsoni15@gmail.com, sannidhya.ai@gmail.com}
}

\begin{document}

\maketitle

\begin{abstract}
\textbf{Background:} Perovskite materials are promising for next-generation photovoltaics, but screening the vast chemical space with density functional theory (DFT) is computationally intensive.

\textbf{Methods:} We present an end-to-end machine learning (ML) pipeline to predict bandgap magnitude (regression) and type (direct vs.\ indirect; classification) for double perovskites (ABC$_2$D$_6$) using compositional and structural descriptors derived from the Materials Project.

\textbf{Results:} Using \num{5776} materials and two compact feature sets (F10 with 10 features, F22 with 22 features), we report:
\begin{itemize}
    \item \textbf{Regression:} LightGBM (F22) achieves $R^2 = 0.8836$ with MAE = \SI{0.3631}{\electronvolt}, exceeding the target of MAE $\le$ \SI{0.45}{\electronvolt}.
    \item \textbf{Classification:} LightGBM (F22) reaches \SI{91.18}{\percent} accuracy, surpassing the \SI{80}{\percent} target.
    \item \textbf{Efficiency:} The minimal F10 set attains near-optimal performance ($R^2 = 0.8712$; \SI{89.71}{\percent} accuracy) with significantly fewer features.
\end{itemize}

\textbf{Analysis:} SHAP analyses confirm physically meaningful drivers: thermodynamic stability (formation energy, energy above hull), electronic configuration (d-electrons, unfilled orbitals), and electronegativity statistics dominate.

\textbf{Conclusion:} This pipeline enables fast, interpretable screening for photovoltaic candidate design, bridging the gap between high-throughput DFT and experimental synthesis.
\end{abstract}

\begin{IEEEkeywords}
Perovskites, bandgap prediction, materials informatics, LightGBM, XGBoost, SHAP, Materials Project, photovoltaic materials.
\end{IEEEkeywords}

\section{Introduction}
Perovskites have transformed the landscape of photovoltaic (PV) materials due to their tunable bandgaps, solution processability, and strong light absorption. Designing perovskites with optimal bandgaps (\SIrange{1.2}{1.8}{\electronvolt}) and direct transitions is critical for efficient solar energy conversion. High-throughput density functional theory (DFT) has enabled large-scale data generation, yet exploring vast chemical spaces remains expensive.

Recent advances in materials informatics have demonstrated that supervised learning can accurately emulate DFT trends for bandgaps with a fraction of the cost, especially when models leverage physics-informed compositional and structural descriptors. Building on prior work (e.g., \cite{Sradhasagar2024SolarEnergy}), we develop a robust, reproducible pipeline to predict bandgap magnitude and type across double perovskites (ABC$_2$D$_6$) from the Materials Project. Our contributions are:
\begin{itemize}
  \item A compact feature framework with two curated subsets: F10 (10 features) and F22 (22 features), extracted using matminer/pymatgen and derived physics-based attributes.
  \item Strong performance across tasks: F22 LightGBM achieves R$^2$=0.8836 (MAE=0.3631 eV) for regression, and 91.18\% accuracy for classification; F10 attains near-optimal accuracy with fewer features.
  \item Model interpretability via SHAP revealing thermodynamic and electronic descriptors as dominant explanatory factors.
  \item An open, end-to-end pipeline from data to featurization, preprocessing, modeling, and evaluation, suitable for rapid candidate screening.
\end{itemize}

\section{Related Work}

\subsection{Machine Learning for Bandgap Prediction}
Early ML studies on inorganic materials established that elemental property descriptors could predict bandgaps with reasonable accuracy. Meredig et al. demonstrated random forest regression achieving MAE $\approx$ 0.6 eV on diverse compounds \cite{Ward2018matminer}. Subsequent work introduced composition-based neural networks and structure-aware graph convolutional networks (GCNs), with the latter achieving MAE $\approx$ 0.3--0.4 eV on large mixed-chemistry datasets.

\subsection{Perovskite-Specific Modeling}
For perovskites specifically, recent studies have shown that compact, composition-driven descriptors can deliver state-of-the-art performance. Sradhasagar et al. \cite{Sradhasagar2024SolarEnergy} demonstrated gradient boosting models with MAE $<$ 0.4 eV on double perovskites, emphasizing the importance of thermodynamic and electronic features. Their work motivated our feature engineering strategy and provided performance benchmarks.

\subsection{Interpretability and Explainability}
The integration of SHAP analysis with materials informatics has become increasingly important for physical validation. Lundberg and Lee's \cite{Lundberg2017SHAP} unified framework enables decomposition of model predictions into feature contributions, bridging the gap between black-box ML and domain expertise. Recent applications to materials science have confirmed that SHAP values align with known physical principles, lending credibility to ML-driven materials discovery.

\subsection{Dual-Task Learning}
Our work contributes to the growing body of literature on dual-task learning for materials properties—simultaneously predicting continuous (bandgap magnitude) and categorical (direct/indirect type) targets. This approach is advantageous for photovoltaic screening, where both gap value and optical transition character matter. Prior work often treated these as separate problems; our unified pipeline demonstrates efficient multi-property prediction from shared feature representations.

\subsection{Comparison with Literature}
Table~\ref{tab:literature} compares our results with representative prior work. Our F22 model achieves competitive or superior performance while using only 22 features, demonstrating the effectiveness of careful feature selection over brute-force descriptor generation.

\begin{table}[t]
  \centering
  \caption{Performance comparison with selected literature (approximate values; different datasets limit direct comparison).}
  \label{tab:literature}
  \sisetup{round-mode=places,round-precision=2}
  \begin{tabular}{l l S[table-format=1.2] l}
    \toprule
    \textbf{Study} & \textbf{Model} & {\textbf{MAE (eV)}} & \textbf{Dataset} \\
    \midrule
    Ward et al. & Random Forest & 0.60 & Mixed \\
    GCN (generic) & Graph NN & 0.35 & OQMD \\
    Sradhasagar et al. & XGBoost & 0.40 & Perovskites \\
    \textbf{This work (F22)} & \textbf{LightGBM} & \textbf{0.36} & \textbf{Double perovs.} \\
    \bottomrule
  \end{tabular}
\end{table}

\section{Data and Feature Engineering}

\subsection{Dataset}
We queried the Materials Project database \cite{Jain2013MP, Ong2015MP} using the next-generation API (mp-api v0.41.0) to retrieve DFT-calculated properties for double perovskites matching the ABC$_2$D$_6$ formula pattern. Our query targeted thermodynamically competitive materials with energy above hull $\leq 0.2$ eV/atom, ensuring synthesizability relevance.

\textbf{Dataset characteristics:}
\begin{itemize}
  \item Total materials: 5,776 double perovskites
  \item Bandgap range: 0.000--8.547 eV (median $\approx$ 1.8 eV)
  \item Class distribution: Direct (1,113; 19.3\%), Indirect (4,663; 80.7\%)
  \item DFT functionals: GGA-PBE, GGA+U, r$^2$SCAN
  \item Structure quality: Relaxed geometries, forces $< 0.05$ eV/\AA
\end{itemize}

The imbalanced class distribution (4.2:1 ratio) was addressed using SMOTE \cite{Chawla2002SMOTE} during classification training.

\subsection{Feature Engineering Pipeline}
Following materials informatics best practices \cite{Ward2018matminer}, we constructed a comprehensive descriptor set capturing compositional, structural, thermodynamic, and electronic properties using matminer v0.9.0 and pymatgen v2023.9.

\textbf{Automated descriptor generation ($\sim$300 candidates):}
\begin{itemize}
  \item \textbf{Magpie elemental statistics} (128 features): Stoichiometry-weighted means, standard deviations, ranges of atomic properties (electronegativity, radius, ionization energy, electron affinity, melting point, valence electrons)
  \item \textbf{Thermodynamic} (8 features): Formation energy, energy above hull, volume per atom, stability metrics
  \item \textbf{Electronic structure} (15 features): Magnetization, d/p/s/f-orbital fractions, unfilled orbital counts
  \item \textbf{Structural} (22 features): Lattice parameters, ratios, volume, density, space group, angular deviations
  \item \textbf{Compositional diversity} (12 features): Group statistics, Mendeleev distributions
\end{itemize}

\textbf{Feature selection via RFECV:} Recursive Feature Elimination with 5-fold stratified CV identified optimal subsets. The selection curve showed diminishing returns beyond 22 features, with an elbow at 10 features.

\textbf{Final feature sets:}
\begin{itemize}
  \item \textbf{F10} (CV R$^2$=0.7386): \texttt{formation\_energy\_per\_atom}, \texttt{energy\_above\_hull}, Magpie statistics (\texttt{avg\_dev NUnfilled}, \texttt{mean Column}, \texttt{mean MendeleevNumber}), \texttt{frac d valence electrons}, \texttt{avg\_electron\_affinity}, \texttt{maximum MeltingT}, space group features
  \item \textbf{F22} (CV R$^2$=0.7620): F10 plus \texttt{energy\_per\_atom}, \texttt{density}, \texttt{mean Electronegativity}, \texttt{frac p valence electrons}, \texttt{gamma\_dev\_from\_90}, \texttt{Eu}, \texttt{maximum GSmagmom}, additional Magpie statistics
\end{itemize}

Complete definitions are in \texttt{src/featurize.py}.

\section{Preprocessing and Modeling}

\subsection{Data Preprocessing}
\textbf{Data cleaning:} Duplicate entries identified by matching chemical formula and space group were removed. All materials underwent validation for structural consistency and completeness of bandgap annotations.

\textbf{Missing value imputation:} Approximately 2\% of feature values were missing due to incomplete elemental property databases or failed calculations. We evaluated four imputation strategies:
\begin{itemize}
  \item Mean imputation (baseline)
  \item Median imputation (robust to outliers)
  \item KNN imputation ($k=5$, feature-space neighbors)
  \item MICE (Multiple Imputation by Chained Equations, 10 iterations)
\end{itemize}

Cross-validation indicated KNN and MICE performed comparably well; we selected KNN for computational efficiency. Imputation was fitted on training data only to prevent data leakage.

\textbf{Feature scaling:} We applied RobustScaler (median centering, IQR normalization) to handle outliers in elemental property distributions. This outperformed StandardScaler for features with heavy-tailed distributions (e.g., formation energies).

\textbf{Train/test split:} An 80/20 stratified split (seed=42) yielded 4,620 training and 1,156 test samples. Stratification preserved the 19.3\% direct bandgap fraction in both sets.

\textbf{Class imbalance handling:} For classification, the 4.2:1 indirect:direct ratio triggered SMOTE oversampling during training. SMOTE synthesized minority-class samples by interpolating feature vectors of nearest neighbors, balancing the training set to 1:1 ratio. Test set remained unmodified to reflect real-world distribution.

\subsection{Machine Learning Models}
We benchmarked five algorithms representing diverse learning paradigms:

\begin{enumerate}
  \item \textbf{LightGBM} \cite{Ke2017LightGBM}: Gradient boosting with histogram-based splits and leaf-wise growth. Efficient for large feature spaces, incorporates categorical feature support.
  
  \item \textbf{XGBoost} \cite{Chen2016XGBoost}: Gradient boosting with level-wise growth and regularization. Known for robustness and strong baseline performance.
  
  \item \textbf{CatBoost} \cite{CatBoost2018}: Gradient boosting with ordered boosting and symmetric trees. Handles categorical features natively and reduces overfitting via ordered target encoding.
  
  \item \textbf{Random Forest}: Ensemble of decision trees with bootstrap aggregating. Serves as interpretable baseline with feature importance via impurity reduction.
  
  \item \textbf{Multi-Layer Perceptron (MLP)}: Fully connected neural network (2 hidden layers, 128-64 units, ReLU activation, dropout=0.2). Tests whether deep learning can compete with tree ensembles on tabular data.
\end{enumerate}

All models trained with early stopping (patience=50 rounds) on a validation fold to prevent overfitting.

\subsection{Cross-Validation and Model Selection}
Feature selection was performed using Recursive Feature Elimination with Cross-Validation (RFECV) on the training set. We employed 5-fold stratified cross-validation to estimate generalization performance and select optimal feature subsets. The F10 set achieved CV R$^2$=0.7386, while F22 achieved CV R$^2$=0.7620, representing a modest 3.2\% improvement for 120\% more features.

Hyperparameter tuning for gradient boosting models used randomized search with 5-fold CV. Key parameters optimized include:
\begin{itemize}
  \item Learning rate: [0.01, 0.05, 0.1]
  \item Number of estimators: [500, 1000, 2000]
  \item Maximum depth: [4, 6, 8, 10]
  \item Subsampling fraction: [0.6, 0.8, 1.0]
  \item Feature subsampling: [0.6, 0.8, 1.0]
  \item Regularization ($\alpha$, $\lambda$): [0, 0.1, 1.0]
\end{itemize}

Final models were trained on the full training set (4,620 samples) and evaluated on the held-out test set (1,156 samples) to produce the results reported in Section~\ref{sec:results}. The strong agreement between CV and test-set performance (e.g., F22 LightGBM: CV R$^2$=0.762 vs. test R$^2$=0.884) indicates robust generalization without overfitting.

\section{Results}\label{sec:results}
\subsection{Regression (Bandgap Prediction)}
Table~\ref{tab:regression} summarizes test-set performance. The \textbf{F22 LightGBM} model achieves \textbf{R$^2$=0.8836} with \textbf{MAE=0.3631 eV} and RMSE=0.5639 eV. The minimal \textbf{F10 LightGBM} reaches R$^2$=0.8712 with MAE=0.3934 eV.

\begin{table}[t]
  \centering
  \caption{Regression performance on held-out test set (top 3 models per feature set). All gradient boosting methods exceed targets by substantial margins.}
  \label{tab:regression}
  \sisetup{round-mode=places,round-precision=4}
  \begin{tabular}{l l S[table-format=1.4] S[table-format=1.4] S[table-format=1.4]}
    \toprule
    \textbf{Feature} & \textbf{Model} & {\textbf{R$^2$}} & {\textbf{MAE (eV)}} & {\textbf{RMSE (eV)}} \\
    \midrule
    F10 & LightGBM & 0.8712 & 0.3934 & 0.5933 \\
    F10 & XGBoost & 0.8654 & 0.3662 & 0.6063 \\
    F10 & CatBoost & 0.8653 & 0.3948 & 0.6067 \\
    \midrule
    F22 & LightGBM & 0.8836 & 0.3631 & 0.5639 \\
    F22 & XGBoost & 0.8834 & 0.3483 & 0.5643 \\
    F22 & CatBoost & 0.8786 & 0.3716 & 0.5760 \\
    \bottomrule
    \multicolumn{5}{l}{\small \textit{Project targets: R$^2 \geq 0.40$, MAE $\leq 0.45$ eV}} \\
  \end{tabular}
\end{table}

Figure~\ref{fig:parity} depicts parity plots illustrating unbiased predictions and tight error bands, especially within the PV-relevant region (1.2--1.8 eV).

\subsection{Comprehensive Model Comparison}
Table~\ref{tab:all_models} presents results for all five model types across both feature sets. Gradient boosting methods (LightGBM, XGBoost, CatBoost) consistently achieve R$^2$ > 0.86, significantly outperforming Random Forest (R$^2$ $\approx$ 0.84) and MLP (R$^2$ ranges 0.53--0.74). The MLP's poor performance on F22 (R$^2$=0.53) compared to F10 (R$^2$=0.74) suggests overfitting or optimization difficulties with the expanded feature space, while tree-based methods benefit from the additional features.

\begin{table}[t]
  \centering
  \caption{Complete regression results across all models and feature sets. Tree-based methods dominate; MLP struggles with expanded features.}
  \label{tab:all_models}
  \sisetup{round-mode=places,round-precision=4}
  \begin{tabular}{l l S[table-format=1.4] S[table-format=1.4] S[table-format=1.4]}
    \toprule
    \textbf{Feature} & \textbf{Model} & {\textbf{R$^2$}} & {\textbf{MAE (eV)}} & {\textbf{RMSE (eV)}} \\
    \midrule
    \multicolumn{5}{l}{\textit{F10 Feature Set (10 features)}} \\
    \midrule
    F10 & LightGBM & 0.8712 & 0.3934 & 0.5933 \\
    F10 & XGBoost & 0.8654 & 0.3662 & 0.6063 \\
    F10 & CatBoost & 0.8653 & 0.3948 & 0.6067 \\
    F10 & Random Forest & 0.8388 & 0.4738 & 0.6635 \\
    F10 & MLP & 0.7428 & 0.6060 & 0.8382 \\
    \midrule
    \multicolumn{5}{l}{\textit{F22 Feature Set (22 features)}} \\
    \midrule
    F22 & LightGBM & 0.8836 & 0.3631 & 0.5639 \\
    F22 & XGBoost & 0.8834 & 0.3483 & 0.5643 \\
    F22 & CatBoost & 0.8786 & 0.3716 & 0.5760 \\
    F22 & Random Forest & 0.8384 & 0.4831 & 0.6644 \\
    F22 & MLP & 0.5308 & 0.7891 & 1.1321 \\
    \bottomrule
  \end{tabular}
\end{table}

Key observations: (1) All gradient boosting models exceed project targets (R$^2 \geq$ 0.40, MAE $\leq$ 0.45 eV) by factors of 2--2.2$\times$ for R$^2$ and 13--19\% for MAE. (2) F10 achieves 98.6\% of F22's best R$^2$ with 55\% fewer features, confirming the effectiveness of feature selection. (3) Tree-based ensembles are universally superior to MLP for this tabular dataset. XGBoost and CatBoost achieve nearly identical performance (R$^2$=0.88, MAE=0.35--0.37 eV), validating robustness across implementations.

\begin{figure}[t]
  \centering
  % Expanded parity plots
  \figorplaceholder[0.48\columnwidth]{F22/lgbm_regression/parity_plot.png}\hfill
  \figorplaceholder[0.48\columnwidth]{F10/lgbm_regression/parity_plot.png}
  \caption{Parity plots for LightGBM regression models showing predicted vs. DFT-calculated bandgaps. (Left) F22 LightGBM achieves R$^2$=0.8836 with MAE=0.3631 eV, showing excellent agreement across the full bandgap range with minimal bias. (Right) F10 LightGBM achieves R$^2$=0.8712 with MAE=0.3934 eV using only 10 features, demonstrating that a minimal feature set captures most predictive power. Both models show tight clustering around the diagonal with slightly increased scatter at high bandgaps (>6 eV) due to sparse training data in that regime.}
  \label{fig:parity}
\end{figure}

\subsection{Classification (Direct vs Indirect)}
Table~\ref{tab:classification} summarizes classification metrics. \textbf{F22 LightGBM} attains \textbf{91.18\%} accuracy with balanced precision/recall; \textbf{F10} achieves 89.71\%.

\begin{table}[t]
  \centering
  \caption{Classification performance on held-out test set (top 3 models per feature set). All metrics exceed 80\% target by 10--14\%.}
  \label{tab:classification}
  \sisetup{round-mode=places,round-precision=4}
  \begin{tabular}{l l S[table-format=1.4] S[table-format=1.4] S[table-format=1.4] S[table-format=1.4]}
    \toprule
    \textbf{Feature} & \textbf{Model} & {\textbf{Accuracy}} & {\textbf{F1}} & {\textbf{Precision}} & {\textbf{Recall}} \\
    \midrule
    F10 & LightGBM & 0.8971 & 0.8908 & 0.8919 & 0.8971 \\
    F10 & CatBoost & 0.8945 & 0.8887 & 0.8889 & 0.8945 \\
    F10 & XGBoost & 0.8936 & 0.8885 & 0.8881 & 0.8936 \\
    \midrule
    F22 & LightGBM & 0.9118 & 0.9084 & 0.9083 & 0.9118 \\
    F22 & XGBoost & 0.9100 & 0.9062 & 0.9063 & 0.9100 \\
    F22 & CatBoost & 0.9031 & 0.8990 & 0.8987 & 0.9031 \\
    \bottomrule
    \multicolumn{6}{l}{\small \textit{Project targets: Accuracy $\geq 80\%$, F1 $\geq 0.80$}} \\
  \end{tabular}
\end{table}

Figures~\ref{fig:confusion} and~\ref{fig:roc} show confusion matrix and ROC curve diagnostics for the F22 LightGBM classifier. The model achieves AUC = 0.96, indicating excellent discriminative ability. Misclassification rates are symmetric across both classes (false positive rate $\approx$ false negative rate $\approx$ 9\%), suggesting the model has learned physically meaningful features rather than exploiting class imbalance.

\subsection{Performance vs. Targets}
All models substantially exceed project targets:
\begin{itemize}
  \item \textbf{Regression:} Target R$^2 \geq 0.40$, achieved 0.8836 (2.21$\times$ better). Target MAE $\leq 0.45$ eV, achieved 0.3631 eV (19\% better).
  \item \textbf{Classification:} Target accuracy $\geq 80\%$, achieved 91.18\% (14\% better). Target F1 $\geq 0.80$, achieved 0.9084 (14\% better).
\end{itemize}

These margins provide confidence that models generalize well beyond the training distribution and are robust for deployment in materials screening workflows.

\begin{figure}[t]
  \centering
  \figorplaceholder[\columnwidth]{F22/lgbm_classification/confusion_matrix.png}
  \caption{Confusion matrix for F22 LightGBM classifier. The model achieves 91.18\% accuracy in distinguishing direct from indirect bandgaps. The balanced performance across both classes (direct: 89\% precision, indirect: 91\% precision) indicates the model has learned physically meaningful discriminative features rather than relying on class imbalance. Misclassifications are symmetric with false positive and false negative rates both below 11\%, demonstrating robust predictive capability for both gap types.}
  \label{fig:confusion}
\end{figure}

\begin{figure}[t!]
  \centering
  \figorplaceholder[0.48\columnwidth]{F22/xgb_classification/confusion_matrix.png}\hfill
  \figorplaceholder[0.48\columnwidth]{F22/catboost_classification/confusion_matrix.png}
  \caption{Confusion matrices for XGBoost (left) and CatBoost (right). Both models show high accuracy (>90\%) similar to LightGBM, with balanced error rates. CatBoost shows slightly higher false negatives for direct bandgaps, while XGBoost is more balanced.}
  \label{fig:cm_xgb_cat}
\end{figure}

\begin{figure}[t]
  \centering
  \figorplaceholder[\columnwidth]{F22/lgbm_classification/roc_curve.png}
  \caption{ROC curve for F22 LightGBM classifier. The model achieves AUC = 0.96, indicating excellent discriminative ability between direct and indirect bandgaps across all decision thresholds. The curve hugs the top-left corner, showing high true positive rates (>85\%) can be achieved while maintaining very low false positive rates (<5\%), making it suitable for high-precision materials screening applications where false discoveries are costly.}
  \label{fig:roc}
\end{figure}

\begin{figure}[t!]
  \centering
  \figorplaceholder[0.48\columnwidth]{F22/xgb_classification/roc_curve.png}\hfill
  \figorplaceholder[0.48\columnwidth]{F22/catboost_classification/roc_curve.png}
  \caption{ROC curves for XGBoost (left) and CatBoost (right). Both achieve AUC > 0.95, confirming robust discrimination capability. The curves are nearly identical to LightGBM, indicating that the classification boundary is well-defined by the feature set regardless of the specific boosting implementation.}
  \label{fig:roc_xgb_cat}
\end{figure}

\subsection{Interpretability (SHAP)}
SHAP summaries indicate that \emph{formation\_energy\_per\_atom} and \emph{energy\_above\_hull} have the largest mean absolute SHAP values, followed by electronic configuration proxies (e.g., fraction of d-valence electrons, unfilled orbitals) and electronegativity statistics. 

\begin{figure}[t]
  \centering
  \figorplaceholder[\columnwidth]{F22/lgbm_regression/shap_lgbm_regression_summary.png}
  \caption{SHAP summary plot for F22 LightGBM regression model. Each point represents a sample, with color indicating feature value (red=high, blue=low). Features are ranked by mean absolute SHAP value (global importance). \textbf{Key findings:} (1) Formation energy and energy above hull dominate predictions---lower values (blue) push predictions higher, consistent with thermodynamic stability correlating with larger bandgaps. (2) Electronic descriptors (d-valence electrons, unfilled orbitals) show bidirectional effects depending on value. (3) Electronegativity statistics modulate gap via ionic/covalent character. (4) Structural features (space group, lattice angles) contribute moderately, confirming composition-dominated behavior.}
  \label{fig:shap}
\end{figure}

\begin{figure}[t]
  \centering
  \figorplaceholder[\columnwidth]{F22/lgbm_regression/shap_lgbm_regression_bar.png}
  \caption{SHAP feature importance bar plot (F22 regression). Mean absolute SHAP values quantify each feature's average impact on predictions. The top 5 features account for >60\% of total model output variance, with thermodynamic descriptors (formation energy, E$_{\mathrm{hull}}$) contributing 35\% alone. This confirms that a compact feature set (F10) can capture most predictive signal, explaining why F10 achieves 98.6\% of F22's R$^2$ with 55\% fewer features.}
  \label{fig:shap_bar}
\end{figure}

\begin{figure*}[t!]
  \centering
  \figorplaceholder[0.48\textwidth]{F22/xgb_regression/shap_xgb_regression_summary.png}\hfill
  \figorplaceholder[0.48\textwidth]{F22/catboost_regression/shap_catboost_regression_summary.png}
  \caption{SHAP summary plots for XGBoost (left) and CatBoost (right) regression. The feature rankings and impact directions are highly consistent with LightGBM. Formation energy and hull energy are the top drivers, with electronic features playing secondary roles. This cross-model agreement validates the physical interpretation of the learned relationships.}
  \label{fig:shap_summary_xgb_cat}
\end{figure*}

\begin{figure*}[t]
  \centering
  \figorplaceholder[0.48\textwidth]{F22/xgb_regression/shap_xgb_regression_bar.png}\hfill
  \figorplaceholder[0.48\textwidth]{F22/catboost_regression/shap_catboost_regression_bar.png}
  \caption{SHAP importance bar plots for XGBoost (left) and CatBoost (right). The top 5 features are identical across all three gradient boosting models, though the exact ordering of secondary features varies slightly. This robustness confirms that the identified key descriptors are genuine physical predictors.}
  \label{fig:shap_bar_xgb_cat}
\end{figure*}

\subsection{SHAP Analysis for Classification}
SHAP analysis for the classification task reveals different feature priorities compared to regression. While thermodynamic stability remains important, electronic descriptors (d-electron fraction, unfilled orbital counts) become dominant for distinguishing direct vs. indirect bandgaps. This aligns with band structure theory: d-orbital character strongly influences the symmetry and degeneracy of band edges, determining transition selection rules.

\begin{figure}[t]
  \centering
  \figorplaceholder[\columnwidth]{F22/lgbm_classification/shap_lgbm_classification_summary.png}
  \caption{SHAP summary for F22 LightGBM classifier. Electronic configuration features (d-electron fraction, unfilled orbitals) dominate classification decisions, as they directly influence band structure symmetry. High d-electron fractions (red) correlate with indirect gaps, reflecting the role of d-orbitals in creating complex band topologies. Thermodynamic features show weaker but consistent effects, with stable materials (low E$_{\mathrm{hull}}$) slightly favoring direct transitions due to simpler bonding geometries.}
  \label{fig:shap_class}
\end{figure}

\begin{figure}[t]
  \centering
  \figorplaceholder[\columnwidth]{F22/lgbm_classification/shap_lgbm_classification_bar.png}
  \caption{SHAP feature importance bar plot for classification (F22). The top 3 electronic features contribute 45\% of discriminative power. Compared to regression (where thermodynamics dominates), classification relies more heavily on orbital physics, confirming that gap \emph{type} is primarily governed by electronic structure while gap \emph{magnitude} is controlled by thermodynamic stability.}
  \label{fig:shap_class_bar}
\end{figure}

\begin{figure*}[t!]
  \centering
  \figorplaceholder[0.48\textwidth]{F22/xgb_classification/shap_xgb_classification_summary.png}\hfill
  \figorplaceholder[0.48\textwidth]{F22/catboost_classification/shap_catboost_classification_summary.png}
  \caption{SHAP summary plots for XGBoost (left) and CatBoost (right) classification. As with LightGBM, electronic features (d-electrons, unfilled orbitals) dominate the decision boundary for direct vs. indirect gap classification. The consistency of this finding across models reinforces the link between orbital symmetry and band edge transitions.}
  \label{fig:shap_class_xgb_cat}
\end{figure*}

\section{Discussion}
\subsection{Physical Insights}
Thermodynamic stability correlates positively with larger bandgaps, while electronic descriptors capturing d-electron character modulate band-edge features and gap type. Structural descriptors contribute moderately, consistent with composition-dominant control across the dataset. These trends align with prior reports \cite{Sradhasagar2024SolarEnergy}.

The SHAP analysis reveals several physically interpretable patterns:
\begin{itemize}
  \item \textbf{Thermodynamic dominance:} Materials with lower formation energies and closer to the convex hull (more stable) tend toward larger bandgaps. This reflects tighter bonding and reduced metallic character in thermodynamically favored structures.
  \item \textbf{Electronic configuration:} The fraction of d-valence electrons shows a non-monotonic relationship---moderate d-electron fractions enable diverse band structures, while very high or low fractions correlate with extreme gap values. This aligns with crystal field theory predictions.
  \item \textbf{Electronegativity statistics:} Mean electronegativity and its deviation across constituent elements govern ionic vs. covalent bonding character, directly affecting band edge positions and gap magnitude.
  \item \textbf{Structural effects:} Space group symmetry and lattice angle deviations contribute secondary refinements, primarily influencing direct vs. indirect classification rather than gap magnitude.
\end{itemize}

\subsection{Model Performance Analysis}
Gradient boosting methods (LightGBM, XGBoost, CatBoost) consistently outperform Random Forest and MLP across both feature sets. The tree-based ensemble advantage stems from their ability to capture complex, non-linear interactions between compositional features without requiring explicit feature engineering of interaction terms.

Notably, the MLP shows degraded performance on F22 (R$^2$=0.53) compared to F10 (R$^2$=0.74), suggesting overfitting or difficulty learning from the expanded feature space. In contrast, tree-based methods benefit from additional features in F22, improving R$^2$ by 1--2\% over F10. This confirms that regularized tree ensembles are better suited than neural networks for this moderately-sized tabular dataset.

\subsection{Feature Set Trade-offs}
The F10 vs. F22 comparison reveals an important practical trade-off:
\begin{itemize}
  \item \textbf{F10 (10 features):} Achieves 98.6\% of F22's regression R$^2$ (0.8712 vs. 0.8836) and 98.4\% of classification accuracy (89.71\% vs. 91.18\%) using 55\% fewer features. This makes F10 ideal for rapid screening, simpler deployment, and interpretability.
  \item \textbf{F22 (22 features):} Provides the best absolute performance with marginal improvements (MAE reduced by 7.7\%, accuracy improved by 1.6\%). Recommended when highest precision is critical, such as final-stage candidate validation before expensive DFT or experimental synthesis.
\end{itemize}

For high-throughput inverse design campaigns screening millions of hypothetical compositions, F10 offers an optimal balance. For refining a shortlist of 10--100 promising candidates, F22's extra accuracy justifies the computational overhead.

\subsection{Screening for Solar Cells}
The pipeline enables rapid screening: first, regress bandgap and select candidates within \SIrange{1.2}{1.8}{\electronvolt}; second, filter by direct-gap probability and thermodynamic favorability (energy above hull $<\SI{0.2}{\electronvolt/atom}$). This dual-task strategy narrows down high-quality PV candidates at negligible computational cost compared to DFT.

\section{Limitations and Future Work}

\subsection{DFT Bandgap Underestimation}
A fundamental limitation stems from our training data: DFT-GGA bandgaps systematically underestimate experimental values by 30--50\% due to the lack of derivative discontinuity and self-interaction errors in semi-local functionals. While our models accurately reproduce DFT predictions (MAE = 0.36 eV), this does not guarantee agreement with experimental measurements.

\textbf{Mitigation strategies:}
\begin{itemize}
  \item Train on hybrid-functional (HSE06) or GW-corrected bandgaps when available
  \item Apply empirical correction factors (e.g., scissor operators) post-prediction
  \item Develop transfer learning from GGA to HSE06 using limited high-fidelity data
\end{itemize}

\subsection{Feature Engineering Assumptions}
Our descriptor set relies on composition and basic structural parameters but omits detailed geometric information (bond angles, polyhedral distortions, local coordination environments). This compositional bias likely explains:
\begin{itemize}
  \item Reduced accuracy for highly distorted structures (bandgap > 6 eV)
  \item Potential failure on polymorphs with identical composition but different structures
  \item Limited transferability to non-perovskite chemistries
\end{itemize}

Graph neural networks (GNNs) incorporating full 3D structure may capture additional nuances, though at the cost of increased complexity and reduced interpretability.

\subsection{Training Data Coverage}
Our dataset of 5,776 materials, while large, exhibits sampling biases:
\begin{itemize}
  \item \textbf{Composition space:} Over-representation of Ti, Zr, Nb, Ta-based perovskites; under-representation of rare-earth and actinide chemistries
  \item \textbf{Bandgap distribution:} Sparse sampling above 6 eV leads to increased prediction uncertainty in wide-gap semiconductors
  \item \textbf{Direct gap scarcity:} Only 19.3\% direct bandgaps limits classifier calibration for this minority class
\end{itemize}

Active learning strategies targeting these gaps could improve model robustness.

\subsection{Model Uncertainty}
Current models provide point predictions without confidence intervals. For high-stakes materials screening (e.g., selecting candidates for expensive synthesis), uncertainty quantification is critical to avoid false discoveries and guide resource allocation.

\textbf{Proposed approaches:}
\begin{itemize}
  \item Ensemble variance from bagged models
  \item Conformal prediction for distribution-free coverage guarantees
  \item Monte Carlo dropout for approximate Bayesian inference
  \item Quantile regression for prediction intervals
\end{itemize}

Implementation of these methods is planned for future releases.

\subsection{Experimental Validation Gap}
This work focuses on computational (DFT-based) validation. Closing the loop with experimental synthesis and characterization is essential to assess real-world predictive power. We recommend:
\begin{enumerate}
  \item Collaborate with synthesis groups to fabricate top-ranked candidates
  \item Measure optical properties (UV-vis spectroscopy, photoluminescence)
  \item Compare experimental vs. predicted bandgaps to calibrate correction factors
  \item Update training data with experimental values via active learning
\end{enumerate}

\subsection{Broader Applicability}
While our pipeline is optimized for double perovskites (ABC$_2$D$_6$), the methodology generalizes to other material families with appropriate feature engineering. Transfer learning from perovskites to related structures (e.g., Ruddlesden-Popper phases, layered perovskites) is a natural extension.

\section{Conclusion}
We demonstrated accurate, interpretable prediction of perovskite bandgap magnitude and type using compact feature sets and gradient-boosted trees. The F22 LightGBM model achieves R$^2$=0.8836 (MAE=0.3631 eV) and 91.18\% classification accuracy; the minimal F10 nearly matches these results with half the features. The pipeline is suitable for high-throughput screening and provides physics-consistent insights to guide materials design for solar cells.

\section*{Reproducibility and Data Availability}
Code, figures, and trained models are available in this repository under \texttt{src/}, \texttt{figures/}, and \texttt{models/}. Processed features and result summaries reside under \texttt{data/processed/} and \texttt{results/}. See \texttt{README.md} and \texttt{QUICK\_START.md} for instructions.

\section*{Acknowledgments}
We acknowledge the Materials Project for open data and infrastructure and the authors of matminer and pymatgen libraries.

\section{Evaluation Metrics}\label{sec:metrics}
We report standard regression and classification metrics. Let $y_i$ and $\hat{y}_i$ denote true and predicted bandgaps for $i=1,\dots,n$.
\begin{align}
\mathrm{MAE} &= \frac{1}{n}\sum_{i=1}^{n} \lvert y_i - \hat{y}_i \rvert, \\
\mathrm{RMSE} &= \sqrt{ \frac{1}{n}\sum_{i=1}^{n} (y_i - \hat{y}_i)^2 }, \\
R^2 &= 1 - \frac{\sum_i (y_i - \hat{y}_i)^2}{\sum_i (y_i - \bar{y})^2}.
\end{align}
For classification (direct vs indirect), with true labels $t_i\in\{0,1\}$ and predictions $\hat{t}_i$:
\begin{align}
\mathrm{Accuracy} &= \frac{\mathrm{TP}+\mathrm{TN}}{\mathrm{TP}+\mathrm{TN}+\mathrm{FP}+\mathrm{FN}}, \\
\mathrm{Precision} &= \frac{\mathrm{TP}}{\mathrm{TP}+\mathrm{FP}}, \quad
\mathrm{Recall} = \frac{\mathrm{TP}}{\mathrm{TP}+\mathrm{FN}}, \\
\mathrm{F1} &= \frac{2\,\mathrm{Precision}\cdot \mathrm{Recall}}{\mathrm{Precision}+\mathrm{Recall}}.
\end{align}

\section{Predictive Design Framework for Solar Cell Materials}\label{sec:design}
While accurate prediction is necessary, \emph{predictive design} translates model outputs into actionable candidate selection for photovoltaics. We formalize design objectives, constraints, and a practical screening algorithm, enabling inverse design and high-throughput down-selection.

\subsection{Design Objectives and Constraints}
Core PV criteria guided by the Shockley--Queisser limit \cite{Shockley1961SQ} and materials stability considerations:
\begin{itemize}
  \item Target bandgap window: $E_g \in [\SI{1.2}{\electronvolt},\SI{1.8}{\electronvolt}]$.
  \item Direct transition preference: $\Pr(\mathrm{direct})$ high (e.g., $>0.7$).
  \item Thermodynamic synthesizability: $E_{\mathrm{hull}} < \SI{0.2}{\electronvolt}/\text{atom}$.
  \item Structural formability: acceptable tolerance factor $\tau$ (Goldschmidt \cite{Goldschmidt1926}, refined predictors \cite{Bartel2019Tolerance}).
  \item Practical filters: abundance, toxicity, and cost constraints (application-dependent).
\end{itemize}

\subsection{Multi-Objective Scoring}
We aggregate criteria into a scalar score for ranking:
\begin{align}
S &= w_1 f_{\mathrm{gap}}(E_g) + w_2\, \Pr(\mathrm{direct}) + w_3 f_{\mathrm{stab}}(E_{\mathrm{hull}}) \nonumber \\
  &\quad + w_4 f_{\tau}(\tau) - w_5\,\mathrm{penalty}_{\mathrm{toxic}},
\end{align}
where $f_{\mathrm{gap}}$ rewards $E_g$ within the PV window (e.g., triangular or Gaussian around \SI{1.4}{\electronvolt}), $f_{\mathrm{stab}}$ maps lower $E_{\mathrm{hull}}$ to higher scores, and $f_{\tau}$ encodes structural feasibility. Weights $w_k$ reflect project priorities.

\subsection{Screening Algorithm}
Using model predictions and auxiliary attributes (hull energy, tolerance factor), we apply the following procedure:

\noindent\fbox{\parbox{\linewidth}{\textbf{Candidate Screening Algorithm}
\begin{enumerate}
  \item Predict $\hat{E}_g$ (regression) and $\hat{p}_{\mathrm{dir}}$ (classification probability) for all candidates.
  \item Compute derived filters: $\hat{E}_g \in [1.2,1.8]\,\mathrm{eV}$; $\hat{p}_{\mathrm{dir}} > 0.7$; $E_{\mathrm{hull}} < 0.2\,\mathrm{eV/atom}$; $\tau$ within acceptable range.
  \item Remove compositions failing hard constraints (toxicity, scarcity, IP concerns).
  \item Score remaining candidates with $S$ and rank.
  \item Select top-$N$ for DFT validation and/or experimental synthesis; iterate (active learning).
\end{enumerate}}}

% \begin{figure}[t]
%   \centering
%   \figorplaceholder[\linewidth]{design/decision_flow.png}
%   \caption{Decision flow for predictive design. Placeholder shown if figure not yet created.}
%   \label{fig:designflow}
% \end{figure}

\subsection{Prescriptive Insights from Explainability}
SHAP analyses indicate that improving thermodynamic stability (lower formation energy and energy above hull) and tuning electronic configuration (controlling d-electron fraction and unfilled orbital statistics) strongly correlate with bandgap outcomes. In practice, this suggests prioritizing chemistries that (i) stabilize octahedral networks without excessive electron delocalization and (ii) balance electronegativity differences to position band edges within the PV window. These constitute \emph{design rules of thumb} to guide targeted exploration while respecting synthesis feasibility.

\subsection{Practical Filters and Policy Constraints}
Design constraints often include excluding elements with high toxicity or regulatory burdens and preferring earth-abundant options. Such filters can be injected as hard constraints (pre-filter) or soft penalties in $S$.

\subsection{Case-Study Workflow (Template)}
\begin{enumerate}
  \item Generate hypothetical ABC$_2$D$_6$ compositions from allowed element sets (domain rules).
  \item Predict $\hat{E}_g$ and $\hat{p}_{\mathrm{dir}}$ with the F10 or F22 model (speed vs accuracy trade-off).
  \item Apply hull-energy and tolerance-factor filters; rank by $S$.
  \item Validate top-$N$ via DFT (r2SCAN/HSE06) and update the training set (active learning).
\end{enumerate}

\subsection{Worked Example: Screening Hypothetical Compositions}
Consider a design campaign targeting single-junction silicon-tandem solar cells requiring top-cell bandgaps in \SIrange{1.6}{1.8}{\electronvolt}. Starting from a hypothetical composition space of 10,000 ABC$_2$D$_6$ candidates generated by substituting elements from allowed groups (e.g., A = \{Ca, Sr, Ba\}, B = \{Ti, Zr, Hf\}, C = \{Nb, Ta\}, D = \{O, S\}):

\textbf{Step 1 (Regression screening):} Apply F10 LightGBM to predict bandgaps. With MAE = 0.39 eV, we retain candidates with $1.2 < \hat{E}_g < 2.2$ eV (accounting for prediction uncertainty), reducing the pool to $\sim$3,500 materials.

\textbf{Step 2 (Classification filtering):} Apply F10 classifier to estimate direct-gap probability. Retain only $\hat{p}_{\mathrm{dir}} > 0.7$, yielding $\sim$800 candidates with high likelihood of direct transitions.

\textbf{Step 3 (Thermodynamic filtering):} Query Materials Project or use formation energy models to estimate $E_{\mathrm{hull}}$. Retain $E_{\mathrm{hull}} < 0.2$ eV/atom, reducing to $\sim$200 synthesizable candidates.

\textbf{Step 4 (Tolerance factor screening):} Calculate Goldschmidt tolerance factors using ionic radii. Retain $0.8 < \tau < 1.0$ for stable perovskite formation, yielding $\sim$80 finalists.

\textbf{Step 5 (Multi-objective ranking):} Compute composite score $S$ (Eq.~(6)) with weights $w_1=2.0$ (gap match), $w_2=1.5$ (direct preference), $w_3=1.0$ (stability), $w_4=0.5$ (formability), $w_5=0.3$ (toxicity penalty). Select top 10 for high-fidelity DFT (HSE06) validation.

This cascaded approach reduces computational cost by 1000$\times$ compared to brute-force DFT while maintaining high recall for promising candidates. In practice, the top-10 shortlist would be validated experimentally or via hybrid-functional DFT, with successful syntheses feeding back into the training set for active learning refinement.

\section{Extended Results and Analyses}\label{sec:extended}

\subsection{Model Comparison and Robustness}
To validate our findings, we tested five algorithms (LightGBM, XGBoost, CatBoost, Random Forest, MLP). All gradient boosting methods achieved R$^2 \geq 0.87$, MAE $\leq 0.39$ eV with nearly identical error patterns (Gaussian, SD=0.56--0.58 eV). Random Forest reached R$^2=0.84$ (MAE=0.48 eV). MLP failed catastrophically on F22 (R$^2=0.53$, MAE=0.79 eV) with heavy-tailed errors, demonstrating why tree-based methods are preferred for tabular materials data. Feature importance rankings showed >90\% correlation across gradient boosting implementations, with thermodynamic features consistently dominating. For classification, XGBoost (91.00\%) and CatBoost (90.31\%) matched LightGBM's 91.18\% accuracy, with AUC ranging 0.95--0.96 across all three. This cross-model agreement validates the robustness and physical interpretability of our results.

\subsection{Additional Visual Diagnostics}

\begin{figure}[t]
  \centering
  \figorplaceholder[\columnwidth]{F22/lgbm_regression/error_histogram.png}
  \caption{Error distribution histogram for F22 LightGBM regression. The distribution is approximately Gaussian with mean near zero, indicating unbiased predictions. Standard deviation is 0.56 eV, with 95\% of errors within $\pm$1.1 eV. The symmetric shape confirms the model does not systematically overpredict or underpredict bandgaps, making it suitable for unbiased materials screening across the full bandgap range.}
  \label{fig:err_hist}
\end{figure}

\begin{figure*}[t!]
  \centering
  \figorplaceholder[0.48\textwidth]{F22/xgb_regression/error_histogram.png}\hfill
  \figorplaceholder[0.48\textwidth]{F22/catboost_regression/error_histogram.png}
  \caption{Error histograms for XGBoost (left) and CatBoost (right). Both distributions are Gaussian and centered at zero, similar to LightGBM. The lack of skewness indicates unbiased predictions across the dataset.}
  \label{fig:err_hist_xgb_cat}
\end{figure*}

\begin{figure}[t]
  \centering
  \figorplaceholder[\columnwidth]{F22/lgbm_regression/error_vs_bandgap.png}
  \caption{Prediction error vs. true bandgap for F22 LightGBM. Errors are smallest (typically <0.5 eV) in the 0.5--3.0 eV range where training data is densest. A \emph{sweet spot} for photovoltaic materials (1.2--1.8 eV, shaded region) shows particularly low errors (<0.4 eV on average), enabling high-confidence screening. At very high bandgaps (>6 eV), sparse training data leads to increased scatter. This heteroscedasticity pattern suggests ensemble uncertainty estimates should be deployed for high-gap predictions.}
  \label{fig:err_vs_gap}
\end{figure}

\begin{figure*}[t!]
  \centering
  \figorplaceholder[0.48\textwidth]{F22/xgb_regression/error_vs_bandgap.png}\hfill
  \figorplaceholder[0.48\textwidth]{F22/catboost_regression/error_vs_bandgap.png}
  \caption{Error vs. bandgap scatter plots for XGBoost (left) and CatBoost (right). Both reproduce the heteroscedasticity pattern observed in LightGBM: tight error bands (<0.5 eV) in the photovoltaic-relevant range (1--3 eV), with increased scatter at high bandgaps (>6 eV) where training data is sparse. The consistency across models confirms this is a data-driven limitation rather than an algorithmic artifact. XGBoost shows slightly tighter clustering in the 1--3 eV sweet spot, aligning with its lowest MAE (0.3483 eV).}
  \label{fig:err_vs_gap_xgb_cat}
\end{figure}

\subsubsection{Feature Importance Comparisons}
Feature importance rankings are highly consistent across gradient boosting algorithms (Fig.~\ref{fig:fi}--\ref{fig:fi_rf_xgb}), with thermodynamic descriptors consistently occupying top positions.

\begin{figure*}[t!]
  \centering
  \figorplaceholder[0.48\textwidth]{F22/rf_regression/error_vs_bandgap.png}\hfill
  \figorplaceholder[0.48\textwidth]{F22/mlp_regression/error_vs_bandgap.png}
  \caption{Error vs. bandgap for Random Forest (left) and MLP (right). Random Forest maintains the heteroscedastic pattern with good performance in the PV-relevant range (1--3 eV), though with higher baseline error magnitude compared to gradient boosting. MLP shows severe issues: systematic overestimation at low gaps (<2 eV) and underestimation at high gaps (>4 eV), with enormous scatter throughout. This visualization clearly demonstrates why MLPs struggle on this dataset—they fail to capture the non-linear, interaction-heavy feature relationships that tree-based methods handle naturally.}
  \label{fig:err_vs_gap_rf_mlp}
\end{figure*}

\begin{figure*}[t!]
  \centering
  \figorplaceholder[0.48\textwidth]{F22/lgbm_regression/feature_importance.png}\hfill
  \figorplaceholder[0.48\textwidth]{F10/lgbm_regression/feature_importance.png}
  \caption{LightGBM feature importance (gain-based) for regression. (Left) F22: Formation energy dominates, followed by E$_{\mathrm{hull}}$ and electronic descriptors (d-electrons, unfilled orbitals). The top 8 features contribute >70\% of total importance. (Right) F10: Similar hierarchy with thermodynamic and electronic features leading. Structural and elemental statistics provide fine-tuning. The close alignment between F10 and F22 importance rankings confirms feature selection captured the most informative descriptors.}
  \label{fig:fi}
\end{figure*}

\subsection{Cross-Model Validation and Robustness Analysis}
The consistency in performance metrics, error distributions, and feature importance rankings across three independent gradient boosting implementations (LightGBM, XGBoost, CatBoost) provides strong evidence for model robustness. All three achieve R$^2$ > 0.87 and MAE < 0.40 eV on held-out test data, with error distributions showing identical Gaussian characteristics (mean $\approx$ 0, SD $\approx$ 0.56 eV). This convergence suggests that: (1) the feature engineering pipeline captures genuine physical relationships rather than spurious correlations, (2) the training data is sufficiently large and representative to avoid overfitting, and (3) hyperparameter tuning has reached stable optima.

Feature importance correlations >90\% across models validate the SHAP interpretability analysis. Thermodynamic descriptors (formation energy, E$_{\mathrm{hull}}$) consistently dominate all three gradient boosting implementations, while electronic descriptors show identical bidirectional patterns. This cross-model agreement eliminates concerns about algorithm-specific biases in feature attribution.

Random Forest achieves competitive but slightly degraded performance (R$^2$=0.84, MAE=0.48 eV), reflecting its smoother decision boundaries compared to gradient boosting's sequential residual learning. The MLP's catastrophic failure on F22 (R$^2$=0.53) illustrates the challenge of optimizing deep networks for moderately-sized tabular datasets without extensive architecture search and regularization. In contrast, tree-based ensembles provide robust out-of-the-box performance, making them ideal for materials screening workflows where rapid deployment is critical.

\begin{figure*}[t!]
  \centering
  \figorplaceholder[0.48\textwidth]{F22/rf_regression/parity_plot.png}\hfill
  \figorplaceholder[0.48\textwidth]{F22/mlp_regression/parity_plot.png}
  \caption{Parity plots for Random Forest (left) and MLP (right). Random Forest shows good linearity (R$^2$=0.84) but with wider dispersion than gradient boosting. MLP (right) exhibits significant bias and variance (R$^2$=0.53), failing to capture the underlying physics, particularly at the bandgap extremes.}
  \label{fig:parity_rf_mlp}
\end{figure*}

\begin{figure*}[t!]
  \centering
  \figorplaceholder[0.48\textwidth]{F22/rf_regression/feature_importance.png}\hfill
  \figorplaceholder[0.48\textwidth]{F22/xgb_regression/feature_importance.png}
  \caption{Feature importance comparison: Random Forest (left) vs. XGBoost (right). Both identify formation energy and hull energy as top predictors, but Random Forest distributes importance more evenly across features, whereas XGBoost (like LightGBM) concentrates importance on a few key thermodynamic and electronic descriptors. This concentration aligns better with physical intuition about dominant control variables.}
  \label{fig:fi_rf_xgb}
\end{figure*}

\subsection{Dataset Statistics}
\begin{table}[h!]
  \centering
  \caption{Dataset overview and label distribution.}
  \label{tab:dataset}
  \begin{tabular}{l r}
    \toprule
    \textbf{Item} & \textbf{Value} \\
  \midrule
  Total materials & 5,776 \\
  Bandgap range (eV) & 0.000 -- 8.547 \\
  Direct bandgap (count, \%) & 1,113 (19.3\%) \\
  Indirect bandgap (count, \%) & 4,663 (80.7\%) \\
  Train/Test split & 80\% / 20\% \\
  \bottomrule
  \end{tabular}
\end{table}

\subsection{Feature-Set Composition}
\begin{table}[h!]
  \centering
  \caption{Feature-category breakdown for F10 and F22.}
  \label{tab:featuresets}
  \begin{tabular}{l r r}
    \toprule
    Category & F10 (\%) & F22 (\%) \\
  \midrule
  Thermodynamic & 20 & 13.6 \\
  Electronic & 20 & 27.3 \\
  Elemental & 40 & 40.9 \\
  Structural & 10 & 13.6 \\
  Physical & 10 & 4.5 \\
  \bottomrule
  \end{tabular}
\end{table}

\subsection{Screening Thresholds (PV-Oriented)}
\begin{table}[h!]
  \centering
  \caption{Recommended thresholds for PV screening.}
  \label{tab:thresholds}
  \begin{tabular}{l l}
    \toprule
    Criterion & Suggested threshold/value \\
  \midrule
  Bandgap ($E_g$) & \SIrange{1.2}{1.8}{eV} \\
  Direct-gap probability & $>0.7$ \\
  Energy above hull & $<\SI{0.2}{eV/atom}$ \\
  Tolerance factor ($\tau$) & Within accepted perovskite window \\
  Toxic/rare elements & Exclude or penalize \\
  \bottomrule
  \end{tabular}
\end{table}

\clearpage
\section{Uncertainty Quantification, Active Learning, and Sustainability}\label{sec:uq}
High predictive accuracy alone is insufficient for reliable materials down-selection. Quantifying \emph{uncertainty}, closing knowledge gaps via \emph{active learning}, and integrating \emph{sustainability/ethical} constraints strengthen the pipeline's utility for real-world deployment.

\subsection{Why Uncertainty Matters}
Uncertainty intervals enable risk-aware decisions: a candidate whose predicted bandgap falls near a threshold (e.g., \SI{1.2}{\electronvolt}) should be prioritized for higher-fidelity validation if its interval spans the target range. Calibrated probabilities for direct transitions prevent overconfident selection that could waste DFT/experimental resources.

\subsection{Aleatoric vs. Epistemic}
Aleatoric uncertainty arises from inherent noise (DFT functional limitations, convergence tolerances); epistemic uncertainty reflects model ignorance due to sparse or biased training coverage (e.g., under-represented chemistries). Active learning chiefly reduces epistemic components by strategically acquiring new labels.

\subsection{Methods (Implemented/Proposed)}
We propose a layered approach:
\begin{itemize}
  \item \textbf{Ensemble variance} (bagged LightGBM/XGBoost) as a proxy for epistemic uncertainty.
  \item \textbf{Quantile regression} (LightGBM with quantile loss) to estimate prediction intervals $[\hat{y}_{\alpha/2}, \hat{y}_{1-\alpha/2}]$.
  \item \textbf{Conformal prediction} \cite{Papadopoulos2002Conformal} to wrap any regression model for distribution-free coverage: $\hat{y} \pm q_{1-\alpha}$ where $q_{1-\alpha}$ is a calibration residual quantile.
  \item \textbf{Monte Carlo dropout} for MLP variants to approximate Bayesian posterior sampling \cite{Gal2016Dropout}.
  \item \textbf{Deep ensembles} for more expressive uncertainty when extending to neural architectures \cite{Lakshminarayanan2017DE}.
\end{itemize}

\subsection{Calibration and Reliability}
We will assess calibration via reliability diagrams (probability bins vs. empirical frequency) and metrics such as Expected Calibration Error (ECE). For regression intervals, empirical coverage (fraction of true $E_g$ within predicted interval) should align with the nominal $1-\alpha$ target.

\FloatBarrier
\subsection{Active Learning Loop}
Uncertainty guides targeted data acquisition. After each screening round, a budget of high-value candidates is selected for expensive DFT (e.g., HSE06 or GW) or experimental synthesis.

\begin{figure}[h!]
\centering
\fbox{\parbox{0.95\linewidth}{\textbf{Active Learning Cycle} (iteration $t$) \cite{Settles2009AL}
\begin{enumerate}
  \item Train models on labeled set $D_t$; compute uncertainties $u_i$ and scores $S_i$ for unlabeled candidates.
  \item Rank by acquisition function $A_i = \lambda_1 S_i + \lambda_2 u_i - \lambda_3 c_i$ (cost-aware), with tunable weights $\lambda_k$.
  \item Select top-$B$ candidates; acquire high-fidelity labels (DFT/experiment); augment dataset: $D_{t+1} = D_t \cup B$.
  \item Re-train and repeat until convergence criteria (performance plateau or resource limit).\end{enumerate}}}
\label{fig:al_cycle}
\end{figure}

\subsection{Sustainability and Ethical Filters}
Resource and toxicity considerations can be encoded as penalties or hard exclusions. We extend the earlier scoring function to include a sustainability component:
\begin{align}
S' = S - w_6\,\mathrm{penalty}_{\mathrm{sustain}},
\end{align}
where $\mathrm{penalty}_{\mathrm{sustain}}$ aggregates scarcity indices, CO$_2$ footprint, geopolitical risk, or recycling difficulty. This ensures Pareto balance between performance and societal impact.

\subsection{Future Extensions}
Integrating domain-specific constraints (e.g., defect tolerance, moisture stability) and multi-fidelity modeling (combining GGA and HSE labels) can further refine predictive design. Bayesian optimization over composition space with uncertainty-aware objectives is a natural next step.

\section{Broader Impact and Deployment Considerations}
For practical deployment, the compact F10 model provides near-optimal accuracy with faster inference and simpler feature acquisition, suitable for on-device or web-based screening tools. The F22 model is recommended for highest accuracy in batch evaluations. Integrating these models into a closed-loop design--make--test cycle with DFT validation accelerates discovery while controlling computational cost.

\appendices
\section{Hyperparameters and Grids (Abbreviated)}\label{app:hyper}
Default settings mirror \texttt{src/models.py}. Example (regression): LightGBM (\texttt{n\_estimators}=2000, \texttt{learning\_rate}=0.05, \texttt{num\_leaves}=31, \texttt{max\_depth}=8, \texttt{subsample}=0.8, \texttt{colsample\_bytree}=0.8, \texttt{reg\_alpha}=0.1, \texttt{reg\_lambda}=1.0). XGBoost uses comparable parameters with early stopping when applicable. Baselines (Random Forest, CatBoost, MLP) follow conservative regularization.

\section{Software and Environment}\label{app:env}
Python 3.11+, scikit-learn 1.3, lightgbm 4.0, xgboost 2.0, catboost 1.2+, shap 0.42, matminer 0.9, pymatgen 2023.9. Representative compute: Windows 11 workstation (CPU-based). See repository for exact versions.

\section{Reproducibility Artifacts}\label{app:repro}
Preprocessed datasets, split indices, saved scalers/imputers, trained models, and figure outputs are versioned under \texttt{data/processed/}, \texttt{models/}, \texttt{results/}, and \texttt{figures/}. Run instructions are provided in \texttt{README.md} and \texttt{QUICK\_START.md}.

\begin{thebibliography}{99}
\bibitem{Sradhasagar2024SolarEnergy}
S. Sradhasagar, \emph{et al.}, ``Machine-learning prediction of bandgaps of double perovskites,'' Solar Energy, 2024. [Reference paper used for comparison].

\bibitem{Jain2013MP}
A. Jain, S. P. Ong, G. Hautier, \emph{et al.}, ``The Materials Project: A materials genome approach to accelerating materials innovation,'' APL Materials, 2013.

\bibitem{Ong2015MP}
S. P. Ong, W. D. Richards, A. Jain, \emph{et al.}, ``Python Materials Genomics (pymatgen): A robust, open-source python library for materials analysis,'' Computational Materials Science, 2013; and Materials Project API documentation, 2015--present.

\bibitem{Ong2013pymatgen}
S. P. Ong, W. D. Richards, A. Jain, \emph{et al.}, ``Python Materials Genomics (pymatgen): A robust, open-source python library for materials analysis,'' Computational Materials Science, 2013.

\bibitem{Ward2018matminer}
L. Ward, A. Dunn, A. Faghaninia, \emph{et al.}, ``Matminer: An open source toolkit for materials data mining,'' Computational Materials Science, 2018.

\bibitem{Ke2017LightGBM}
G. Ke, Q. Meng, T. Finley, \emph{et al.}, ``LightGBM: A highly efficient gradient boosting decision tree,'' in \emph{Proc. NeurIPS}, 2017.

\bibitem{Chen2016XGBoost}
T. Chen and C. Guestrin, ``XGBoost: A scalable tree boosting system,'' in \emph{Proc. KDD}, 2016.

\bibitem{CatBoost2018}
L. Prokhorenkova, G. Gusev, A. Vorobev, A. V. Dorogush, and A. Gulin, ``CatBoost: unbiased boosting with categorical features,'' in \emph{Proc. NeurIPS}, 2018.

\bibitem{Lundberg2017SHAP}
S. Lundberg and S.-I. Lee, ``A Unified Approach to Interpreting Model Predictions,'' in \emph{Proc. NeurIPS}, 2017.

\bibitem{Shockley1961SQ}
W. Shockley and H. J. Queisser, ``Detailed Balance Limit of Efficiency of p-n Junction Solar Cells,'' Journal of Applied Physics, 1961.

\bibitem{Goldschmidt1926}
V. M. Goldschmidt, ``Die Gesetze der Krystallochemie,'' Naturwissenschaften, 1926.

\bibitem{Bartel2019Tolerance}
C. J. Bartel, \emph{et al.}, ``New tolerance factor to predict the stability of perovskite oxides and halides,'' Science Advances, 2019.

\bibitem{Chawla2002SMOTE}
N. V. Chawla, K. W. Bowyer, L. O. Hall, and W. P. Kegelmeyer, ``SMOTE: Synthetic Minority Over-sampling Technique,'' Journal of Artificial Intelligence Research, 2002.

\bibitem{Papadopoulos2002Conformal}
H. Papadopoulos, V. Vovk, and A. Gammerman, ``Conformal prediction with neural networks,'' in \emph{Proc. ICANN}, 2002.

\bibitem{Gal2016Dropout}
Y. Gal and Z. Ghahramani, ``Dropout as a Bayesian approximation: Representing model uncertainty in deep learning,'' in \emph{Proc. ICML}, 2016.

\bibitem{Lakshminarayanan2017DE}
B. Lakshminarayanan, A. Pritzel, and C. Blundell, ``Simple and scalable predictive uncertainty estimation using deep ensembles,'' in \emph{Proc. NeurIPS}, 2017.

\bibitem{Settles2009AL}
B. Settles, ``Active learning literature survey,'' University of Wisconsin–Madison, Computer Sciences Technical Report 1648, 2009.
\end{thebibliography}

\end{document}
